\section{Nothing On That Route Asks}
\label{sec:ch12-nothing-on-that-route-asks}

\index{entities@\code{\_entities}!what the specification says about it|(}%
Apollo's subgraph specification defines the route in one line, and
chapter~\ref{ch:first-subgraph} printed it:

\begin{graphqlsdl}[fontsize=\footnotesize]
_entities(representations: [_Any!]!): [_Entity]!
\end{graphqlsdl}

\index{entities@\code{\_entities}!what the specification validates}%
The page goes on to say what a subgraph library must validate, and the list
is entirely about the representation coming in: that it carries
\code{\_\_typename}, and that it carries every field named in the type's
\code{@key} field set~\autocite{apollosubgraphspec}. Those are questions
about whether a representation is well formed. Neither of them is a question
about who sent it.

\index{entities@\code{\_entities}!the words that are not in the specification}%
The words \enquote{authentication} and \enquote{authorization} do not appear
on that page at all. It lists the definitions of \code{@authenticated},
\code{@requiresScopes} and \code{@policy} among the directives a subgraph
schema may carry, and says nothing about applying any of them to the entity
route, nothing about who enforces them, and nothing about what a subgraph
should do with a representation from a caller it does not recognize. I went
looking for that sentence in the same way chapter~\ref{ch:representation-order}
went looking for one about checking the answer. It is not there.

\index{entities@\code{\_entities}!the route has no identity to check}%
Which is consistent rather than careless, once you look at what the route is
for. Everything about its signature says machine. The argument is
\code{\_Any}, a scalar that exists so representations can be passed through
untyped. The return is a union with no common field. There is no user in that
design, because the caller the design has in mind is a router that has
already decided the request is allowed.

\subsection{Apollo Says So Outright}

\index{subgraph!not to be reachable by clients}%
The assumption is not left implicit. Apollo's own security guidance states it
as a rule and names this field as the reason:
\enquote{As a best practice for supergraphs, only the router should query
individual subgraphs directly}, because \enquote{if this field is exposed
publicly, any client can circumvent internal resolver logic and fetch any
entity data by mimicking the router}~\autocite{apollographsecurity}.

\index{subgraph!the network position is a control}%
That is the whole architecture in two sentences, and it is a real control. A
subgraph nobody can route a packet to cannot be asked anything. If your three
services are on a private network and the router is the only thing with a
public address, the requests in
section~\ref{sec:ch12-three-doors-one-column} are requests nobody can send.

\index{subgraph!a network position is not a guard}%
I would still not leave it there, for a reason that has nothing to do with
distrusting the network team. A network position is a property of a
deployment, and the guard it substitutes for is a property of the code. They
live in different repositories, change on different schedules, and are
reviewed by different people. The code moves to a new environment; the
position does not come with it. Anything already inside the network is inside
it, which covers a compromised sidecar, a debugging proxy somebody left up, a
developer machine on the VPN, and the next service that gets deployed into
the same subnet. And there is no test that fails when the position stops
holding.

\index{subgraph!both, and the argument is not about which}%
So both, and I do not think that is a compromise. The network position is
what stops a stranger reaching the port. The guard in the service is what
means it did not matter that they got there. They fail in different ways and
they are not substitutes for one another, which is the only reason to have
two of anything.

\index{authorization!nobody bylined writes about this route}%
I looked for somebody to hand this argument to and did not find one. Two
named engineers make the general case, and I take both: David Walter at
Apollo, that subgraphs \enquote{should only be accessible from your router,
never directly from the public internet}~\autocite{walter2026subgraphs}, and
Tom Houl\acu{e} at Grafbase, a competing implementation, warning that
\enquote{if subgraphs accept unauthenticated traffic, attackers could bypass
security by directly accessing subgraph
endpoints}~\autocite{houle2026federation}. Neither
of them writes about \code{\_entities} in particular. The one article I found
that does is signed with a company name and no person's, and I am not willing
to hang a security claim on writing nobody put their name to. So the specific
claim in this chapter, that this route in this stack walks past a guard on a
root field, is mine. What stands behind it is the measurements above rather
than an authority, and if that is not enough for you, the requests are short
and your own subgraph is right there.

\subsection{The Router Does Hide It}

\index{router!hides the entity route from clients}%
One thing chapter~\ref{ch:enter-the-router} established is worth repeating
here, because it cuts the other way. Introspect the router rather than a
subgraph and the route is not there:

\begin{minted}[fontsize=\footnotesize,breakanywhere]{json}
{"data":{"__schema":{"queryType":{"fields":[{"name":"node"},
 {"name":"nodes"},{"name":"sessions"},{"name":"speakers"},
 {"name":"ratingCount"}]}}}}
\end{minted}

\index{router!hiding is not guarding}%
\code{\_entities} and \code{\_service} are absent from the composed schema,
which is what the specification asks for and what the router does. A client
talking to the router cannot name the field, cannot construct the query, and
cannot see it in the schema.

That is the graph's public face and it is correct. It is also not a guard on
the service, which is still listening on its own port with the same three
doors it had before. Hiding a field from one client is not the same as
refusing it to another.
\index{entities@\code{\_entities}!what the specification says about it|)}%
